\documentclass[11pt]{article}
\usepackage[T1]{fontenc}
\usepackage{lmodern}
\usepackage[margin=1in]{geometry}
\usepackage{microtype}
\usepackage{amsmath,amssymb,amsthm}
\usepackage[hidelinks]{hyperref}

\newtheorem{theorem}{Theorem}
\newtheorem{lemma}{Lemma}
\newcommand{\gt}{\gamma_t}
\newcommand{\repobase}{https://github.com/burgerge/wowii255-total-domination}
\newcommand{\archivebase}{https://github.com/PalomarArchive/burgerge--wowii255-total-domination--ccc373ad5fc0}
\newcommand{\catalogue}{https://www.uhd.edu/academics/sciences/mathematics-statistics/wowii/index.html}

\title{Minimum Counterexample Orders and Sharp Coefficients\\
       for Two Total Domination Conjectures}
\author{Chenyang Zhang}
\date{27 November 2026}
\hypersetup{
  pdftitle={Minimum Counterexample Orders and Sharp Coefficients for Two Total Domination Conjectures},
  pdfauthor={Chenyang Zhang}
}

\begin{document}
\maketitle

\begin{abstract}
Conjectures 255 and 256 of \emph{Written on the Wall II} propose
lower bounds on the total domination number of a connected graph in
terms of its center and the neighborhood of its minimum-degree vertices,
respectively. We prove that the minimum counterexample orders are
\(11\) and \(14\). In each inequality the largest valid universal
coefficient is \(1\), and the normalized ratio has infimum \(1\)
even over connected bipartite graphs. The proofs use neighborhood
covering arguments and two explicit graph families. Lean
formalizations accompany the minimum-order and sharpness results.
\end{abstract}

\section{The conjectures and results}

Let \(G\) be a finite connected simple graph of order \(n\ge 2\).
Write \(N(v)\) for the open neighborhood of \(v\), and
\(N(S)=\bigcup_{v\in S}N(v)\) for any vertex set \(S\).
A set \(D\) is total dominating when \(N(D)=V(G)\). In particular,
the vertices in \(D\) must also have neighbors in \(D\).
The minimum size of such a set is the total domination number
\(\gt(G)\). For an edge \(uv\) write \(B(uv)=N(u)\cup N(v)\) and set
\[
 M(G)=\max_{uv\in E(G)}|B(uv)|.
\]
The set \(B(uv)\) includes both endpoints of the edge. Let \(C(G)\)
be the set of vertices of minimum eccentricity, and let \(A(G)\)
be the set of vertices of minimum degree.

The two conjectures in DeLaVi\~na's catalogue~\cite{wowii} take the
following form after multiplication by the positive denominator \(M(G)\).
\begin{align}
 \gt(G)M(G)&\ge 2|C(G)|,   \tag{255}\label{eq:255}\\
 \gt(G)M(G)&\ge 2|N(A(G))|. \tag{256}\label{eq:256}
\end{align}
Blair published an eleven-vertex counterexample to
\eqref{eq:255} and a fifteen-vertex counterexample to
\eqref{eq:256}~\cite{blair255,blair256}. Earlier public notes
give counterexample families for both conjectures~\cite{priorfamilies}.
For \eqref{eq:256}, the catalogue records Tony Newton's fourteen-vertex
witness, a proof of its minimality, and an extremal-coefficient result
submitted on 21 November 2026~\cite{newton256}. The unrestricted
optimum \(1\) for \eqref{eq:256} also appears in the 19 November
manuscript of the Graffiti.pc Conjecture 256 repository
contributors~\cite{prior256}, using graphs with triangles.

We give complete proofs of the two minimum orders and two sharp
coefficients. In particular, the construction for \eqref{eq:256}
shows that the coefficient stays sharp on connected bipartite
graphs. The formalization includes both minimum orders and the
bipartite strengthening for \eqref{eq:256}.

\begin{theorem}\label{thm:min}
The minimum orders of counterexamples to \eqref{eq:255} and
\eqref{eq:256} are \(11\) and \(14\), respectively. Every graph
of order \(2\le n\le 10\) satisfies the stronger inequality
\(\gt(G)M(G)\ge 2n\).
\end{theorem}

\begin{theorem}\label{thm:sharp}
Let \(T(G)\) be either \(C(G)\) or \(N(A(G))\), and let \(\alpha\) be real.
The inequality \(\gt(G)M(G)\ge\alpha|T(G)|\) holds for every
finite connected simple graph of order at least two precisely
when \(\alpha\le 1\). For each choice of \(T\), the ratio
\(\gt(G)M(G)/|T(G)|\) is strictly greater than \(1\) on every
such graph and tends to \(1\) along a family of connected
bipartite graphs.
\end{theorem}

\section{Neighborhood covering bounds}

Write \(g=\gt(G)\) and \(M=M(G)\), and take a minimum total dominating
set \(D\). The induced graph \(G[D]\) has no isolated vertex.
If a collection of its edges covers \(D\), their sets \(B(uv)\)
cover \(V(G)\). In particular, if \(g=2\), then \(M=n\).

\begin{lemma}\label{lem:degree}
If \(M<n\), then \(\deg(u)\le M-2\) for every vertex \(u\).
\end{lemma}

\begin{proof}
Suppose that \(Q=\{u\}\cup N(u)\) has size at least \(M\).
For every \(v\in N(u)\) we have \(Q\subseteq B(uv)\), so the
definition of \(M\) forces \(Q=B(uv)\). Thus all neighbors of
vertices of \(Q\) belong to \(Q\). Connectedness gives
\(Q=V(G)\), contradicting \(M<n\).
\end{proof}

If \(gM<2n\), then \(g\ge 2\) gives \(M<n\), and the preceding
observations give \(g\ge 3\). The sets \(N(v)\) for \(v\in D\)
cover the graph, so Lemma~\ref{lem:degree} implies
\begin{equation}\label{eq:cover}
 n\le g(M-2),\qquad n+2g\le gM<2n,\qquad 2g<n.
\end{equation}
If \(g=3\), we can label \(D=\{a,b,c\}\) with \(ab,bc\in E(G)\).
The two sets \(B(ab),B(bc)\) cover the graph and both contain \(D\).
Consequently
\begin{equation}\label{eq:path}
 n+3\le 2M.
\end{equation}
In the equality case, the two sets each have size \(M\) and
intersect exactly in \(D\). In particular \(N(b)=\{a,c\}\),
as \(N(b)\) is contained in their intersection.

\section{Minimum counterexample orders}

\subsection{An eleven-vertex witness for Conjecture 255}

For \(s\ge 2\) let \(H_s\) have vertices \(a,b,c\) and
\(x_i,y_i\) for \(1\le i\le s\). Its edges are
\[
 ab,\quad bc,\quad ax_i,\quad x_i y_i,\quad y_i c,\quad
 x_i x_j \quad(1\le i<j\le s).
\]
The path arrangement without the clique on the \(x_i\) occurs in the
catalogue's example for the neighboring Conjecture 254~\cite{wowii254}.
The family for \(s\ge4\) appears in~\cite{priorfamilies}. An equivalent
description in~\cite{evan255} subdivides each edge joining a vertex of
the clique \(\{a,x_1,\ldots,x_s\}\) to the hub \(c\). The case \(s=4\)
gives Blair's example~\cite{blair255}.

All distances in \(H_s\) are at most two. The \(x_i\) form a clique,
the \(y_i\) share the neighbor \(c\), and for \(i\ne j\) the pair
\(x_i,y_j\) is linked by \(x_j\). The pairs \(a,y_i\), \(b,x_i\),
\(b,y_i\), and \(c,x_i\) use \(x_i,a,c,y_i\), respectively, as
common neighbors. The pair \(a,c\) uses \(b\). As the largest
degree is \(s+1<2s+2\), every vertex has eccentricity exactly
two and \(|C(H_s)|=2s+3\).

With \(X=\{x_1,\ldots,x_s\}\) and \(Y=\{y_1,\ldots,y_s\}\), the
open neighborhoods of \(a,b,c\) are \(X\cup\{b\}\), \(\{a,c\}\),
and \(Y\cup\{b\}\). Those of \(x_i,y_i\) are
\((X\setminus\{x_i\})\cup\{a,y_i\}\) and \(\{x_i,c\}\).
Their unions at the six edge types all have size \(s+3\), giving
\(M(H_s)=s+3\). The set \(\{a,b,c\}\) totally dominates \(H_s\),
and \(M(H_s)<2s+3\) excludes a total dominating pair. Thus
\[
 \gt(H_s)=3,\qquad \frac{\gt(H_s)M(H_s)}{|C(H_s)|}
   =\frac{3(s+3)}{2s+3}.
\]
For \(s=4\) this gives \(3\cdot7=21<22\). The examples with \(s\ge
4\) all violate \eqref{eq:255}. Equation~\eqref{eq:path} also
shows that \(3/2\) is the infimum of this ratio among graphs with
\(g=3\), attained in the limit \(s\to\infty\).

\subsection{The lower bound for Conjecture 255}

Suppose \(n\le10\) and \(gM<2n\). By \eqref{eq:cover}, we have
\(g\in\{3,4\}\). For \(g=3\), \eqref{eq:path} and \(3M<2n\)
give \(3(n+3)<4n\), so \(n=10\). Then \(M\ge7\), contradicting
\(3M<20\). For \(g=4\), the inequality \(4M<2n\le20\) gives
\(M\le4\), and \(n\le4(M-2)\le8\), contradicting \(2g<n\).
This proves the stronger lower bound \(gM\ge2n\) and establishes
the claimed minimum order for \eqref{eq:255}.

\subsection{A fourteen-vertex witness for Conjecture 256}

We use Newton's witness as recorded in the catalogue~\cite{newton256}.
Take a path \(a,b,c\) and a disjoint cycle \(z_0,z_1,\ldots,z_{10}\)
in cyclic order. Add the edges \(bz_0\), \(az_i\) for \(1\le i\le5\),
and \(cz_i\) for \(6\le i\le10\). The resulting graph is connected.
The vertices \(a,c\) have degree \(6\), and every other vertex has
degree \(3\). Every vertex has a neighbor of degree \(3\), so
\(N(A(G))=V(G)\). Each edge joins vertices whose degrees sum to
at most \(9\). At \(ab\) the neighborhoods of its endpoints
are disjoint and have sizes \(6\) and \(3\), so \(M=9\).
The path vertices form a total dominating set, and \(M<14\)
excludes a total dominating pair. Hence \(g=3\), and
\(gM=27<28=2|N(A(G))|\).

\subsection{The lower bound for Conjecture 256}

Suppose \(n\le 13\) and \(gM<2|N(A(G))|\le 2n\).
By \eqref{eq:cover}, \(g\) is one of \(3,4,5,6\).
We use the following observation for the middle two cases.
In a graph on at least four vertices with no isolated vertices,
if any two edges intersect then the graph is a star. Indeed, take
two edges \(xy,xz\). Any edge avoiding \(x\) must be \(yz\).
If \(yz\) is present, an edge incident to a vertex outside the
triangle cannot intersect all its edges. The absence of isolated
vertices excludes this, so every edge is incident to \(x\).

\paragraph{\(g=3\).}
By \eqref{eq:path}, \(n+3\le2M\), and by assumption \(3M<2n\le26\).
These bounds imply \(M\in\{7,8\}\). If \(M=7\), then \(n=11\),
and if \(M=8\), then \(n=13\). In both cases equality holds
in \eqref{eq:path}. With \(D=\{a,b,c\}\) as above, we have
\(N(b)=\{a,c\}\) and \(|B(ab)|=|B(bc)|=M\). Thus
\(\deg(a),\deg(c)\ge M-2\ge5\), while \(\delta(G)\le\deg(b)=2\).
Neither neighbor of \(b\) has minimum degree, so \(b\notin N(A(G))\).
For both pairs \((n,M)\), this gives
\(3M\ge2(n-1)\ge2|N(A(G))|\), a contradiction.

\paragraph{\(g=4\).}
Two disjoint edges of \(G[D]\) would cover \(D\) and give \(n\le2M\),
contradicting \(4M<2n\). Thus \(G[D]\) is a four-vertex star.
The three sets \(B(uv)\) for its edges cover \(V(G)\) and each
contains \(D\), giving \(n\le3M-8\). The bound \(4M<2n\le26\)
gives \(M\le6\). If \(M\le5\), then \(n\le7\), contradicting
\(2g<n\). If \(M=6\), then \(n\le10\) and \(4M<2n\) forces
\(n>12\), again a contradiction.

\paragraph{\(g=5\).}
From \(5M<2n\le26\) and \(10<n\le5(M-2)\) we obtain \(M=5\)
and \(n=13\). Every vertex has degree at most three, so \(G[D]\)
cannot be a five-vertex star. Take two disjoint edges and an edge
joining the fifth vertex to a neighbor in \(D\). The three edges
cover \(D\). Two of them intersect, and the two corresponding
sets \(B(uv)\) have at least the three vertices of their union of
endpoints in common. It follows that \(n\le3M-3=12\).

\paragraph{\(g=6\).}
The bound \(6M<2n\le26\) gives \(M\le4\), so \(n\le6(M-2)\le12\)
contradicts \(2g<n\).

This excludes counterexamples on at most thirteen vertices and
completes the proof of Theorem~\ref{thm:min}.

\section{The sharp coefficients}

For every \(v\in V(G)\), choose a neighbor \(w\). Then \(v\in N(w)\)
and \(v\notin N(v)\), so \(\deg(v)+1\le|B(vw)|\le M\).
For either target set \(T(G)\), the neighborhood cover by \(D\) yields
\begin{equation}\label{eq:strict}
 |T(G)|\le n\le\sum_{v\in D}\deg(v)\le g(M-1)<gM.
\end{equation}
Both target sets are nonempty. This proves that coefficient \(1\)
is valid and that equality in the normalized ratio is impossible.
The two families below prove sharpness.

\subsection{The center}

Take integers \(r\ge2\) and \(m\ge1\). Let
\(I=\{a_1,\ldots,a_r,b_1,\ldots,b_r\}\).
Make a bipartite incidence graph with point vertices \(I\times I\).
For each label \(i\in I\), make \(m\) distinct block vertices for
\[
 S_i=\{(x,y)\in I\times I: x=i\ \text{or}\ y=i\}.
\]
Add two large blocks whose point sets are
\((I\times I)\setminus\{(a_1,a_1)\}\) and
\((I\times I)\setminus\{(b_1,a_1)\}\). Adjacency is point
membership in a block.

Two points share a large block except for the two distinct omitted
points, which share a small block of label \(a_1\).
Two small blocks of labels \(i,j\) contain \((i,j)\).
A small block of label \(i\) and either large block contain
\((i,b_1)\), and both large blocks contain \((b_1,b_1)\).
It follows that distances within either part are at most two.
Each block has a neighbor, so distances between the parts are at
most three. Each point misses a small block of a label absent from
its coordinates, and each block misses at least one point.
Nonadjacent vertices in opposite parts have distance at least three.
The graph is therefore connected and every vertex has eccentricity
three. In particular \(|C(G)|=4r^2+2+2rm\).

The \(r\) points \((a_i,b_i)\) meet all small blocks. Together with
the two large blocks they form a total dominating set. Thus \(g\le
r+2\). The degrees of points and blocks are bounded by \(2m+2\)
and \(4r^2\), respectively, so \(M\le2m+4r^2+2\).
Set \(Q=4r^2+2\) and \(m=rQ\). For \(r\ge2\),
\[
 (r+2)(2m+Q)=(2r^2+5r+2)Q\le(2r+6)m.
\]
Hence \(1<gM/|C(G)|\le1+3/r\). Taking \(r\to\infty\) proves
sharpness for the center.

\subsection{The neighborhood of the minimum-degree vertices}\label{sec:minimum-degree-family}

Choose integers \(k,s,t\ge2\). In each of \(k\) disjoint groups,
take \(s\) original left vertices and \(st\) ordinary right
vertices, with all edges between these two sets. Partition the
ordinary right vertices into \(t\) batches of size \(s\).
Add one left vertex per batch and join it to the right vertices
in that batch. Finally add a right vertex \(z\), adjacent to every
left vertex in all groups. The graph is bipartite and connected.

The degrees of the original left vertices, batch left vertices,
ordinary right vertices, and \(z\) are \(st+1,s+1,s+1,k(s+t)\),
respectively. The vertices of minimum degree are the batch left
vertices and the ordinary right vertices. Every vertex has a
neighbor of minimum degree, giving
\[
 |N(A(G))|=k(st+s+t)+1.
\]
Every total dominating set contains at least one left vertex per
group to cover its ordinary right vertices, and at least one right
vertex to cover left vertices. One original left vertex per group
together with \(z\) suffices, so \(g=k+1\).

The neighborhoods of adjacent vertices in a bipartite graph are
disjoint. The largest degrees in the left and right parts occur at
an original left vertex and \(z\), respectively, and those
vertices are adjacent. Hence \(M=st+1+k(s+t)\).
For example, \((k,s,t)=(2,4,5)\) gives a bipartite graph of order
\(59\) with \(gM=117<118=2|N(A(G))|\).
For \(s=t=k^2\) the ratio is
\[
 \frac{\gt(G)M(G)}{|N(A(G))|}
  =\frac{(k+1)(k^4+2k^3+1)}{k^5+2k^3+1}\longrightarrow 1
       \qquad(k\to\infty).
\]
Together with \eqref{eq:strict}, this proves the remaining part of
Theorem~\ref{thm:sharp}.

\section{Formal verification and availability}

The Lean development proves both minimum orders, both sharp
coefficients, and the sharp coefficient for Conjecture 256 when
restricted to bipartite graphs. The formal proof for Conjecture 256
uses the equal-parameter subfamily from
Section~\ref{sec:minimum-degree-family} and a rational upper bound
tending to one. The public repository provides a theorem index
and independent integer checks of the finite witnesses~\cite{code}.

The Palomar registrations use fixed source snapshots checked by the Lean
kernel and the independent NanoDa kernel.
\begin{description}
 \item[Sharp coefficients.]
 \href{\archivebase/tree/f1808b7d12e18b22b003b11acf0b3668f24d0bf4}{\texttt{PALOMAR-2026-11-27-000006}}.
 \item[Minimum orders.]
 \href{\archivebase/tree/fd8be52bd12973419e94c1a7471aa7dc508a9c3f}{\texttt{PALOMAR-2026-11-27-000011}}.
\end{description}
The joint source release is available at
\href{https://doi.org/10.5281/zenodo.24567987}{DOI 10.5281/zenodo.24567987}.

\paragraph{Use of generative AI.}
OpenAI Codex was used in developing constructions and proofs,
formalizing and checking the results in Lean, and preparing the manuscript.

\begingroup
\small
\begin{thebibliography}{9}

\bibitem{wowii}
E.~DeLaVi\~na,
\emph{Written on the Wall II: Conjectures of Graffiti.pc},
Conjectures 255 and 256, University of Houston--Downtown.
\href{\catalogue}{Online catalogue}.

\bibitem{wowii254}
E.~DeLaVi\~na,
\emph{Written on the Wall II}, Conjecture 254 and its
\href{https://www.uhd.edu/academics/sciences/mathematics-statistics/wowii/files/cedalm254.jpg}{counterexample diagram}.

\bibitem{blair255}
D.~Blair, \emph{An 11-vertex counterexample to Graffiti.pc
Conjecture 255}, public proof and Lean certificate, 25 November 2026.
\href{https://github.com/derekblair/math-proofs/blob/af6904cc3fa537b5c4dcea9561b24b69177c6ee1/research/lean-verified-2026-11-25/graffiti-255/PROOF255.md}{Fixed source}.

\bibitem{blair256}
D.~Blair, \emph{A 15-vertex counterexample to Graffiti.pc
Conjecture 256}, public proof and Lean certificate, 25 November 2026.
\href{https://github.com/derekblair/math-proofs/blob/af6904cc3fa537b5c4dcea9561b24b69177c6ee1/research/lean-verified-2026-11-25/graffiti-256/PROOF.md}{Fixed source}.

\bibitem{priorfamilies}
\texttt{f8bsd9yqgp-maker}, \emph{Solved Conjectures and Solved Math},
24 November 2026. Notes on
\href{https://github.com/f8bsd9yqgp-maker/Solved-conjectures-/blob/9a49179c1217e4e260e24b9d5d506432d1b328eb/proofs/graffiti-pc-255/README.md}{Conjecture 255}
and
\href{https://github.com/f8bsd9yqgp-maker/Solved-conjectures-/blob/9a49179c1217e4e260e24b9d5d506432d1b328eb/proofs/graffiti-pc-256/README.md}{Conjecture 256}.

\bibitem{evan255}
EvanDashRust, \emph{Graffiti.pc Conjecture 255 is false},
Conjecture-Killers, 25 November 2026.
\href{https://github.com/EvanDashRust/Conjecture-Killers/blob/1c777ae198b5a9fcddc69968232d0eddd4138f80/Results/Graffiti255/README.md}{Fixed source}.

\bibitem{newton256}
T.~Newton, results submitted to the \emph{Written on the Wall II}
catalogue, 21 November 2026. See E.~DeLaVi\~na's
\href{https://www.uhd.edu/academics/sciences/mathematics-statistics/wowii/files/comments.json}{comment on Conjecture 256}
and the
\href{https://www.uhd.edu/academics/sciences/mathematics-statistics/wowii/files/ce256.png}{counterexample diagram}.

\bibitem{prior256}
Graffiti.pc Conjecture 256 repository contributors,
\emph{Counterexamples to Graffiti.pc Conjecture 256 and a sharp total
domination bound}, 19 November 2026.
\href{https://github.com/srikant07-dev/graffiti-pc-conjecture-256/blob/314f74a234703f3860327586054c9fa9e08e45e7/paper.pdf}{Fixed manuscript}.

\bibitem{code}
C.~Zhang, \emph{Written on the Wall II 255 and 256:
Lean formalizations of smallest counterexamples and sharp coefficients},
2026. \href{\repobase/tree/67878cbe235830fc174bb13c0d03b77d104742f5}{Formal proofs
and independent verification code}.

\end{thebibliography}
\endgroup
\end{document}
